\addchap{Úvod}

V oblasti malých a středních servisních firem je evidence zakázek a vystavování faktur často řešena pomocí jednoduchých tabulek, papírových formulářů nebo nejednotných softwarových nástrojů. Tento přístup s sebou přináší zvýšenou administrativní zátěž, riziko chyb při ručním přepisování údajů a omezený přehled o aktuálním stavu jednotlivých zakázek. Zaměstnanci tak tráví značné množství času rutinními úkony, které by bylo možné efektivně automatizovat.

Na trhu existuje řada webových i desktopových aplikací určených pro fakturaci a správu zakázek. Tyto systémy však bývají často navrženy jako univerzální řešení pro široké spektrum uživatelů, což vede k jejich složitosti a nutnosti přizpůsobovat se cizímu pracovnímu postupu. Pro menší servisní firmy proto nemusí být tato řešení ideální, neboť ne vždy odpovídají jejich specifickým potřebám a zavedeným pracovním procesům.

Cílem této práce je analyzovat současná softwarová řešení pro evidenci servisních zakázek a fakturaci, identifikovat jejich přínosy a nedostatky z pohledu konkrétní servisní firmy a na základě této analýzy navrhnout, implementovat a nasadit vlastní webovou aplikaci přizpůsobenou reálnému provozu firmy. Navržené řešení má umožnit jednotnou evidenci zakázek, provedených prací a použitých materiálů a jejich následné využití při automatizované tvorbě nabídek a faktur.

Aplikace je vyvíjena a zveřejněna jako open-source software pod licencí MIT, takže si ji může zdarma nasadit a přizpůsobit i jiná malá firma.

Teoretická část nejprve popisuje průběh servisní zakázky a fakturační procesy v malé firmě a shrnuje existující softwarová řešení. Dále uvádí zásady návrhu přehledného a efektivního uživatelského rozhraní, principy databázového návrhu a přehled použitých technologií.

Praktická část začíná analýzou požadavků firmy a popisem workflow zakázky. Následuje návrh architektury, databázového modelu a uživatelského rozhraní aplikace, ve kterém je návrh porovnán s doporučeními z teoretické části. Další dvě kapitoly popisují implementaci systému pro evidenci zakázek a automatizované tvorby nabídek a faktur z uložených dat. Poslední kapitola popisuje testování, nasazení aplikace do provozu a vyhodnocení zpětné vazby od zaměstnanců.

